\subsection{Urban commuting and public transit}
\label{ch:extensions}

The freight application uses goods-market clearing and mobile labor to map
transport costs into welfare. Urban applications instead use residence and
workplace choices. Because both settings share the transport block, the
extension changes the spatial residuals and welfare projection while leaving
route and mode aggregation intact.

\subsubsection{Urban implementation and verification}

The urban application changes the spatial model, not the transport
technology. In the freight model, locations produce and consume traded goods.
In the urban model, workers choose residences and workplaces and travel
between them. These choices produce different equilibrium residuals and a
different welfare row. For any given residence--workplace flow, however,
commuters face the same recursive routes, modal logsum, and congestion
feedback as shipments in the freight model.

Section \ref{ch:equilibrium} defines the residence--workplace model, and
Subsection \ref{ch:ift} shows how its Jacobian and welfare row enter the shared
transport derivative. The implementation stores those model-specific
quantities in \code{UrbanEngine}. The route, modal, congestion, and primitive
pass-through maps remain in the common transport modules.

The one-mode regression oracle independently codes the Allen--Arkolakis
urban derivative \citep{AA_2022restud}. The shared engine must reproduce its
state response and welfare derivative below \(10^{-10}\), including the
\(1+\theta_U\gamma\) congestion denominator. To see the role of this
denominator, consider one road mode with edge-local congestion elasticity
\(\gamma\). A lower primitive edge cost attracts additional traffic, which
raises the realized cost. Route demand responds with curvature
\(\theta_U\), so the congestion feedback attenuates the primitive shock by
the factor \(1/(1+\theta_U\gamma)\). Reproducing this factor verifies that
the shared transport operator has the same local feedback as the original
one-mode urban system.

The multimodal urban example adds road and transit to the same model and
compares the analytical IFT with an independently solved nonlinear system
under \code{NC}, \code{NT}, \code{F}, \code{FM}, and \code{FR}. The tests allow
residence and workplace margins to differ and include both road and transit
modes under edge and endpoint-terminal congestion. They also verify the
one-mode, unique-route, no-congestion, and zero-terminal limits, as well as
invariance to node and mode permutations.

Run the self-contained example with:
\begin{lstlisting}[language=bash]
julia --project=. bin/tnw.jl validate \
  examples/urban_multimodal/config.toml
julia --project=. bin/tnw.jl decompose \
  examples/urban_multimodal/config.toml
\end{lstlisting}

\subsubsection{Endpoint-terminal congestion}

Edge congestion assigns traffic feedback to a particular edge-mode. Terminal
congestion aggregates affected flows using a terminal and applies the cost
response to incident edge-mode pairs. The terminal identifiers are explicit
input fields; geographic co-location is not enough.

In the freight extension, rail traffic loads the terminals at both endpoints
with elasticity \(0.096\), estimated by
\citet{FuchsWong2026Multimodal}. The main paper result excludes terminal
congestion. Sensitivity to terminal congestion is therefore an extension
diagnostic, not a headline parameter sensitivity.

\subsubsection{Seattle multimodal candidate}

\warningbox{The Seattle module is a candidate empirical application. Its
multimodal urban derivative and accounting checks are implemented, but route
ridership, generalized modal costs, Seattle-specific substitution, and
congestion parameters require additional evidence before the results can be
used as an empirical appraisal.}

The Seattle workflow combines the published 217-location Allen--Arkolakis
urban geography \citep{AA_2022restud} with 2017 LODES
residence--workplace commuting \citep{CensusLODES} and 2017 ACS residential
commute-mode shares \citep{CensusACS2017B08301}. A hash-pinned June 2017 King
County GTFS feed supplies the transit network and schedules
\citep{KingCountyGTFS2017}. The Metro system evaluation enters only as an
external comparison \citep{KingCountyMetro2017Evaluation}.

The baseline uses \(\eta=1.099\), transferred from the freight application in
\citet{FuchsWong2026Multimodal}. It is not estimated from Seattle data.
Baseline modal shifters reproduce the constructed shares; \(\eta\) controls the
local response after an improvement.
The workflow therefore reports sensitivity at \(0.75\), \(0.90\), \(1.099\),
\(1.25\), and \(1.40\).

GTFS describes stops, trips, shapes, and schedules. It does not provide
passenger counts. ACS supplies residential mode shares, not route-level
ridership. The builder routes a common commuter population to obtain a
model-consistent road and transit baseline, records fallback and coverage, and
uses the Metro table only as an external comparison.

\subsubsection{Transit policies}

The workflow evaluates bus, Link light rail and streetcar, and ferry policies
separately and together. A link policy lowers primitive generalized cost on
one directed transit edge-mode. Reciprocal pairs are aggregated only when both
directions exist.

A named route corridor changes several edge-mode costs at once. Let \(b\)
denote its set of edge-mode pairs, \(w_{e,m}\) the relative cost-reduction
weight assigned to pair \((e,m)\), and \(E_{e,m}\) its marginal welfare
elasticity. Linearity of the first-order derivative gives the elasticity of
the bundled reduction:
\[
 E_b=\sum_{(e,m)\in b}w_{e,m}E_{e,m}.
\]
Shared segments change the aggregate mode cost on that segment. The result is
not interpreted as a shock exclusive to one operator's vehicles.

Maps distinguish road, bus, Link light rail, Seattle Streetcar, and ferry by
color. Line width represents the primitive-cost welfare gain, not ridership.
Subway or light-rail vision maps that include proposed future lines are not a
baseline description of the 2017 network.

\subsubsection{Remaining empirical work}

The package computes and verifies the multimodal urban derivative, but a
publication-quality Seattle transit application still requires observed
passenger counts by route or segment and mode-specific measures of generalized
cost and service quality. It also needs a Seattle-specific estimate or external
calibration of \(\eta\), evidence on crowding or terminal-congestion
elasticities, and validation against commuting and ridership outcomes that were
not used to construct the baseline.
These requirements concern measurement and calibration, not the IFT method.
Until they are met, the Seattle module tests the implemented
residence--workplace derivative and multimodal transport operators but does
not provide an empirical transit appraisal. These empirical boundaries also
govern how results from the Seattle module should be described. Section
\ref{ch:conclusion} explains how they enter interpretation and reporting.
